% 写作提示：综述与题目直接相关的文献、理论或标准，并说明材料从何而来。
\section{文献、理论与材料}
\label{sec:context}

\subsection{相关研究或理论基础}
不同课程会把本节写成文献综述、理论框架或基础知识。示例保留三条可以并列的路径，并用\cite{gbt7714-2015,lamport1978,vaswani2017,knuth1986,hushidong2025}演示国标、期刊、会议、图书和网络文献的引用。正式写作时，改成与题目直接相关的文献、标准或案例。

\begin{table}[htbp]
  \centering
  \caption{常见分析路径及其适用情形}
  \label{tab:paths}
  \begin{tabular}{@{}l >{\raggedright\arraybackslash}p{0.36\textwidth} >{\raggedright\arraybackslash}p{0.40\textwidth}@{}}
    \toprule
    \heiti 路径 & \heiti 主要做法 & \heiti 更适用的任务 \\
    \midrule
    实验测量 & 控制条件，并记录可以重复的观测 & 对象能够测量，结果需要对照 \\
    文本细读 & 围绕材料做分层解释 & 论点依赖原文、档案或作品 \\
    案例比较 & 用同一标准并列多个对象 & 需要解释差异及其条件 \\
    \bottomrule
  \end{tabular}
\end{table}

\tabref{tab:paths}里的路径可以改成理论流派、实验方案、场地样本或设计方向。表格的作用是把差异放到同一组标准下。

\subsection{研究对象或资料来源}
示例没有使用外部数据。方案 A、B、C 的分项分数写在附录\tabref{tab:raw}，正文只引用汇总结果。若课程要求说明样本、文本、档案或实验材料，应在这里交代来源、筛选条件，以及不能使用的部分。

\subsection{概念界定与评价标准}
三项指标都按 0 到 1 记录，越高越好。完整性表示对象是否被覆盖，准确性表示判断是否有材料支持，表达表示结果是否清楚、可否复核。三项等权，计算方式见\eqref{eq:score}。
